\subsection{Growth bounds for nonsplit indefinite binary spaces}
\begin{theorem}\label{th:realgrowth}
For a fixed nonsplit indefinite binary parent and either orientation, for every fixed \(\epsilon>0\),
\[
 X^{3/2-\epsilon}\ll_{L,\epsilon}\sum_{m\le X}a_H(m)\ll_L X^2/\log X\qquad(X\ge3).
\]
The abscissa lies in \([3/2,2]\). Along real \(s=2+x\) the series is \(O_L(1+\log(1/x))\) as \(x\downarrow0\), so it cannot have a pole at two.
\end{theorem}
\begin{proof}
\emph{Upper bound.}\quad Realize the space as a line over the real quadratic field \(F\) generated by the parent rotation \(\varepsilon\) of Theorem~\ref{th:inertreal}, acting as multiplication by a norm-one unit. Its eigenvalues are distinct and of absolute values \(\lambda,\lambda^{-1}\), where \(\lambda>1\). Counting orbits of this rotation on raw subgroups of index at most \(X\) gives an upper bound for the number of intrinsic classes. Those of length at least \(J\) contribute \(O(X^2/J)\) by the binary HNF divisor sum. 
Every shorter orbit consists of points fixed by some \(\varepsilon^j\), \(1\le j\le J\). For a lattice \(M\) fixed by that power put \(I=O_F M\subset O_F L\). There are \(O_L(X\log(2X))\) possible \(I\): their index in \(O_F L\) is at most \([O_F L:L]X\) and invertible ideal arithmetic reduces this count to integral ideals, bounded coefficientwise by the divisor function by prime splitting. Since \(\varepsilon^j\) is not rational, \(R_j=\mathbb Z[\varepsilon^j]\) is a quadratic order and \(f_j=[O_F:R_j]\ll_L \lambda^{O(j)}\) by its discriminant \((\varepsilon^j-\varepsilon^{-j})^2\). We have \(f_jI\subset M\subset I\). The number of intervening additive subgroups is bounded by a fixed power of \(1+f_j\), for instance by counting all HNFs of indices at most \(f_j^2\). 
Thus the total short-orbit contribution is \(O_L(e^{c_LJ}X\log(2X))\). Take \(J\) a sufficiently small positive multiple of \(\log X\).

\medskip
\noindent\emph{Lower bound.}\quad Let \(\epsilon>0\) and take \(0<\delta<\min(\epsilon/8,1/16)\). Choose a fixed modulus \(M_0\) divisible by eight and the discriminant of \(F\), so primes \(p\equiv1\bmod M_0\) split in \(F\). Enlarge \(M_0\) to include all bad primes. For large \(y\) consider primes \(q\) in
\([y^{1/2-2\delta},y^{1/2-\delta}]\). There are \(\gg_{L,\delta}y/\log y\) distinct split primes \(p\in(y/2,y]\) for which some such \(q\) divides \(p-1\). In fact sum the primes in residue class \(1\bmod M_0q\) over these \(q\). By the Bombieri--Vinogradov theorem \cite{IK}, applied at both \(y\) and \(y/2\) and using partial summation, the total error from the expected prime counts is \(O_{L,A}(y/(\log y)^A)\) for any fixed \(A\). All moduli are within its range for large \(y\), since \(M_0\) is fixed. The main term is \((\operatorname{Li}(y)-\operatorname{Li}(y/2))\sum_q1/\varphi(M_0q)\), of order \(y/\log y\) by Mertens' theorem over the indicated interval of prime \(q\). An individual \(p\) is counted only \(O_\delta(1)\) times since the product of its distinct divisors \(q\) is at most \(p-1\). This proves the assertion.

Put \(z=\lceil y^{1/2+2\delta}\rceil\). For each of these split primes, \(p-1=q h\) with \(q,h\le z\), so \(p-1\) divides \(\operatorname{lcm}(1,\ldots,z)^2\). Choose products \(m\) of exactly \(k=\lfloor y^{1/2+3\delta}\rfloor\) distinct such primes, and put \(X_y=y^k\). At each \(p\mid m\) there are \(p-1\) primitive normals of nonzero norm, by the split binary finite-field formula. Every isometry between the resulting regular kernels stabilizes the restored parent. The projective parent rotation group is cyclic. Its local projective order divides \(p-1\): a proper split binary rotation preserves both isotropic lines and is diagonal with entries \(t,t^{-1}\) over \(\mathbb F_p\). The simultaneous orbit length is therefore at most \(\operatorname{lcm}(1,\ldots,z)^2\). Consequently either orientation counts at least
\[
 \frac{(y/2-1)^k}{2\operatorname{lcm}(1,\ldots,z)^2}
   =X_y^{1-o(1)}
\]
classes at each of these indices. The ordinary count is always at least half the proper one. Here even \(\log\operatorname{lcm}(1,\ldots,z)\le z\log z=o(k\log y)\) suffices. Unique factorization distinguishes the indices. Their number is at least
\(\binom{\lfloor c_{L,\delta}y/\log y\rfloor}{k}
 \ge X_y^{1/2-3\delta-o(1)}\).
They all lie below \(X_y\), and \(\log X_{y+1}/\log X_y\to1\) for integer \(y\). Monotonicity extends the lower bound to all real \(X\) with the stated loss \(\epsilon\). Partial summation proves the assertions about the abscissa and real boundary growth. Positivity and Landau's theorem show that, if the abscissa is two, this real boundary point is singular despite the absence of a pole there.
\end{proof}


